\section{Problem Statement}
\label{sec:problem_statement}

Despite their growing utility, the widespread deployment of UAVs in cooperative and autonomous scenarios is fundamentally hindered by persistent challenges related to security, scalability, and operational efficiency. The decentralized and dynamic nature of FANETs makes them particularly susceptible to a range of cyber-physical attacks. Malicious actors can exploit the open wireless medium to launch attacks such as data spoofing, message replay, and denial-of-service (DoS), which can compromise mission integrity, lead to loss of assets, or even endanger public safety \cite{luo2023escm}. Centralized communication architectures, while simpler to manage, introduce a single point of failure and are often impractical for large-scale, distributed FANET operations. Consequently, there is a critical need for a decentralized trust and security framework that can operate effectively in highly dynamic and potentially adversarial environments.

Furthermore, the resource-constrained nature of many UAV platforms imposes strict limitations on computational power, energy consumption, and communication bandwidth. Traditional blockchain implementations, particularly those employing computationally intensive consensus mechanisms like PoW, are ill-suited for such environments, as they introduce significant latency and energy overhead \cite{hossain2024blockchain}. This creates a fundamental tension between the desire for robust, blockchain-enforced security and the necessity for lightweight, energy-efficient protocols that do not compromise the UAV's primary mission objectives. The research problem, therefore, is to design and evaluate a blockchain-based communication framework for UAVs that simultaneously provides strong security guarantees while satisfying the stringent performance requirements of FANETs in terms of latency, throughput, scalability, and energy efficiency. This necessitates a holistic approach that re-examines the blockchain protocol stack, from data encoding and consensus to security and network management, to create a solution specifically tailored for the unique characteristics of UAV communication networks.
